% Problem record op_4a4434b5cc6e85e6. This TeX file is derived from
% database/problems_json/op_4a4434b5cc6e85e6.json by scripts/sync-tex.mjs; edit the JSON record
% and rerun the script rather than editing this file.
\section{Gaussian entanglement of formation beyond bisymmetry}
\label{sec:op_4a4434b5cc6e85e6}

\paragraph{Problem.}
Does Gaussian entanglement of formation equal unrestricted entanglement of formation for every non-bisymmetric multimode Gaussian state? Let $\rho_{AB}$ be Gaussian, with integers $n_A,n_B\geq1$ and $n_A+n_B\geq3$. Assume finite total mean photon number. For $S(\tau):=-\operatorname{Tr}(\tau\log_2\tau)$, define
\begin{equation}
E_F(\rho_{AB}):=\inf_{\mu}\int
S(\operatorname{Tr}_B|\psi\rangle\langle\psi|)\,\mu(d\psi).
\label{eq:p26-entanglement-of-formation}
\end{equation}
The probability measures in Eq.~\eqref{eq:p26-entanglement-of-formation} range over normalized pure vectors with barycentre $\rho_{AB}$. Non-Gaussian vectors are allowed.

The Gaussian restriction is equivalently
\begin{equation}
E_F^{\mathrm G}(\rho_{AB}):=
\inf_{\substack{V_p\preceq V\\V_p\ \mathrm{pure\ Gaussian}}}E(V_p),
\label{eq:p26-gaussian-entanglement-of-formation}
\end{equation}
where $V$ is the covariance of $\rho_{AB}$ and $E(V_p)$ is the entropy of either reduced state of the pure Gaussian component. A covariance is bisymmetric when it is invariant under all mode permutations within Alice's block and, independently, within Bob's block. Is $E_F(\rho_{AB})=E_F^{\mathrm G}(\rho_{AB})$ outside this family, with the quantities defined by Eqs.~\eqref{eq:p26-entanglement-of-formation} and \eqref{eq:p26-gaussian-entanglement-of-formation}?

\subsection*{Status}

Unsolved

\subsection*{Source}

Adesso's preprint proves Gaussian optimality for all two-mode and bisymmetric multimode states and leaves the generic nonsymmetric multimode case open \sourcecite{ref:p26-adesso}{Ade26}. The equivalent universal Gaussian-discord formulation below follows from purification duality, rather than being a second independent problem \sourcecite{ref:p26-koashi}{KW04}, \sourcecite{ref:p26-discord}{AD10}.

\subsection*{Progress}

\begin{itemize}
  \item Wolf, Giedke, Kr\"uger, Werner, and Cirac introduced the Gaussian
  convex-roof restriction and derived the covariance-matrix optimization in
  Eq.~\eqref{eq:p26-gaussian-entanglement-of-formation}.  Their analysis does
  not establish equality with the unrestricted roof in
  Eq.~\eqref{eq:p26-entanglement-of-formation}
  \sourcecite{ref:p26-wolf}{WGKWC04}.

  \item Adesso's arXiv preprint establishes $E_F=E_F^{\mathrm G}$ for all two-mode Gaussian states, then extends it to bisymmetric multimode states by localization to a two-mode core. Theorem 1 and Eqs. (18)--(22) state these results; Supplement S2.4 includes continuous pure-state ensembles. Generic nonsymmetric multimode states remain open \sourcecite{ref:p26-adesso}{Ade26}.

  \item For the discord formulation, Adesso--Datta Eqs. (3)--(4) solve the Gaussian measurement optimization for two modes \sourcecite{ref:p26-discord}{AD10}. Pirandola et al., Eqs. (12)--(14), prove unrestricted measurement optimality for states obtained by sending one half of a two-mode squeezed vacuum through a phase-insensitive Gaussian channel. Heterodyne detection on the other half attains the minimum. Their further Gaussian-unitary decompositions cover additional states, but do not settle the universal multimode problem \sourcecite{ref:p26-optimal}{PSBCL14}.
\end{itemize}

\subsection*{References}

\begin{enumerate}[label={},leftmargin=4.8em,labelsep=0.5em]
  \footnotesize
  \item[\textup{[WGKWC04]}]\label{ref:p26-wolf}
  M. M. Wolf, G. Giedke, O. Kr\"uger, R. F. Werner, and J. I. Cirac,
  ``Gaussian Entanglement of Formation,'' \emph{Physical Review A}
  \textbf{69}, 052320 (2004).
  \href{https://doi.org/10.1103/PhysRevA.69.052320}{doi:10.1103/PhysRevA.69.052320};
  \href{https://arxiv.org/abs/quant-ph/0306177}{arXiv:quant-ph/0306177}.

  \item[\textup{[Ade26]}]\label{ref:p26-adesso}
  G. Adesso, ``Optimality of Gaussian Entanglement of Formation,''
  arXiv:2608.01909v2 (2026).
  \href{https://doi.org/10.48550/arXiv.2608.01909}{doi:10.48550/arXiv.2608.01909};
  \href{https://arxiv.org/abs/2608.01909v2}{arXiv:2608.01909v2}.

  \item[\textup{[KW04]}]\label{ref:p26-koashi}
  M. Koashi and A. Winter, "Monogamy of Quantum Entanglement and Other Correlations," \emph{Physical Review A} \textbf{69}, 022309 (2004). \href{https://doi.org/10.1103/PhysRevA.69.022309}{doi:10.1103/PhysRevA.69.022309}; \href{https://arxiv.org/abs/quant-ph/0310037}{arXiv:quant-ph/0310037}.

  \item[\textup{[AD10]}]\label{ref:p26-discord}
  G. Adesso and A. Datta, "Quantum versus Classical Correlations in Gaussian States," \emph{Physical Review Letters} \textbf{105}, 030501 (2010). \href{https://doi.org/10.1103/PhysRevLett.105.030501}{doi:10.1103/PhysRevLett.105.030501}; \href{https://arxiv.org/abs/1003.4979}{arXiv:1003.4979}.

  \item[\textup{[PSBCL14]}]\label{ref:p26-optimal}
  S. Pirandola, G. Spedalieri, S. L. Braunstein, N. J. Cerf, and S. Lloyd, "Optimality of Gaussian Discord," \emph{Physical Review Letters} \textbf{113}, 140405 (2014). \href{https://doi.org/10.1103/PhysRevLett.113.140405}{doi:10.1103/PhysRevLett.113.140405}; \href{https://arxiv.org/abs/1309.2215}{arXiv:1309.2215}.
\end{enumerate}

\subsection*{Comment}

The generic multimode equality remains open. The resolved two-mode case is recorded separately as \texttt{01M26KH5NF45HBHXTPG8ZBD75H}; its resolution is currently an arXiv preprint.

An equivalent universal question asks whether Gaussian measurements attain one-sided discord for every finite-energy bipartite Gaussian state $\omega_{AX}$. For a POVM $\{M_x\}$ on $X$, set $p_x=\operatorname{Tr}[(I_A\otimes M_x)\omega_{AX}]$ and $\omega_{A|x}=\operatorname{Tr}_X[(I_A\otimes M_x)\omega_{AX}]/p_x$ for $p_x>0$. Define
\begin{equation}
D_X(\omega_{AX}):=S(\omega_X)-S(\omega_{AX})
+\inf_M\int p_x S(\omega_{A|x})\,dx .
\label{eq:p26-discord}
\end{equation}
All POVMs are allowed in Eq.~\eqref{eq:p26-discord}; integrals include discrete sums. Define $D_X^{\mathrm G}$ by restricting to general-dyne measurements, including joint multimode seeds and homodyne limits.

For a finite-energy Gaussian purification $\Psi_{ABX}$, the ensemble-measurement correspondence gives
\begin{equation}
D_X^{\mathrm G}(\Psi_{AX})-D_X(\Psi_{AX})
=E_F^{\mathrm G}(\Psi_{AB})-E_F(\Psi_{AB}).
\label{eq:p26-duality}
\end{equation}
The ordinary identity follows from Koashi--Winter Theorem 1, Eq. (2); Adesso--Datta discuss its Gaussian version after Eq. (7) \sourcecite{ref:p26-koashi}{KW04}, \sourcecite{ref:p26-discord}{AD10}. For the Gaussian restriction, every pure covariance $V_p\preceq V$ in Eq.~\eqref{eq:p26-gaussian-entanglement-of-formation} can be prepared by a rank-one Gaussian measurement on a Gaussian purification. Singular boundary cases use homodyne limits; mixed measurement seeds cannot improve the minimum over their pure refinements. The unrestricted correspondence permits probability-measure ensembles.

Every finite-energy Gaussian state has a finite-mode, finite-energy Gaussian purification. Thus Eq.~\eqref{eq:p26-duality} identifies the two universal optimality questions. This equivalence, together with the resolved two-mode and bisymmetric cases, is why the discord formulation is consolidated here.

\subsection*{Field}

Quantum Resource Theory

\subsection*{Topic}

Gaussian quantum information; Entanglement measures

\subsection*{ID}

\texttt{op\_4a4434b5cc6e85e6}
